% Modified on 6-4-05.

%\setcounter{chapter}{11}

\chapter{Jalan Acak}\label{chp 12}  

\section{Jalan Acak dalam Ruang Euklides}\label{sec 12.1}

\par
Dalam beberapa bab terakhir, kita telah mempelajari jumlah peubah acak dengan tujuan
mendeskripsikan fungsi distribusi dan fungsi kepadatan dari jumlah tersebut.  Dalam bab ini, kita akan melihat
jumlah peubah acak diskret dari sudut pandang yang berbeda.  Kita akan menelaah
sifat-sifat yang dapat dikaitkan dengan barisan jumlah parsial, seperti banyaknya
perubahan tanda pada barisan ini, banyaknya suku dalam barisan yang sama dengan 0, dan
nilai harapan suku maksimum dalam barisan tersebut.
\par
Kita mulai dengan definisi berikut.
\begin{definition} 
 Misalkan $\{X_k\}_{k = 1}^\infty$ adalah barisan peubah acak diskret yang saling bebas dan
berdistribusi identik.  Untuk setiap bilangan bulat positif $n$, kita gunakan $S_n$ untuk menyatakan jumlah $X_1 + X_2 +
\cdots + X_n$.  Barisan $\{S_n\}_{n = 1}^\infty$ disebut \emx {jalan
acak.\index{jalan acak}}  Jika daerah nilai bersama dari $X_k$ adalah ${\mathbf R}^m$, maka kita mengatakan bahwa
$\{S_n\}$ adalah jalan acak dalam ${\mathbf R}^m$.
\end{definition} 
\par
Kita memandang barisan $X_k$ sebagai hasil-hasil percobaan yang saling bebas.  Karena
$X_k$ saling bebas, probabilitas suatu barisan hasil tertentu (yang berhingga) dapat
diperoleh dengan mengalikan probabilitas bahwa setiap $X_k$ mengambil nilai yang ditentukan dalam
barisan tersebut.  Tentu saja, probabilitas individual ini diberikan oleh distribusi bersama
$X_k$.  Biasanya kita akan tertarik mencari probabilitas kejadian yang melibatkan
barisan $S_n$ yang terkait.  Kejadian semacam itu dapat dideskripsikan dalam $X_k$, sehingga
probabilitasnya dapat dihitung menggunakan gagasan di atas.
\par
Ada beberapa cara untuk memvisualisasikan jalan acak.  Kita dapat membayangkan sebuah partikel
ditempatkan di titik asal dalam ${\mathbf R}^m$ pada waktu $n = 0$.  Jumlah $S_n$ menyatakan posisi
partikel pada akhir detik ke-$n$.  Jadi, dalam selang waktu $[n-1, n]$, partikel
bergerak (atau melompat) dari posisi $S_{n-1}$ ke $S_{n}$.  Vektor yang menyatakan gerakan ini adalah
$S_n - S_{n-1}$, yang sama dengan $X_n$.  Artinya, dalam suatu jalan acak, lompatan-lompatannya
saling bebas dan berdistribusi identik.  Sebagai contoh, jika $m = 1$, kita dapat membayangkan sebuah 
partikel pada garis real yang mulai dari titik asal dan, pada akhir setiap detik, melompat satu
satuan ke kanan atau ke kiri, dengan probabilitas yang diberikan oleh distribusi $X_k$.  Jika
$m = 2$, prosesnya dapat divisualisasikan berlangsung di sebuah kota yang jalan-jalannya membentuk blok-blok
persegi.  Seseorang mulai dari satu sudut (yaitu, di persimpangan dua jalan) dan berjalan ke
salah satu dari empat arah yang mungkin menurut distribusi $X_k$.  Jika $m = 3$,
kita dapat membayangkan berada di sebuah kerangka panjat, tempat kita bebas bergerak ke salah satu dari enam arah
(kiri, kanan, maju, mundur, atas, dan bawah).  Sekali lagi, probabilitas gerakan-gerakan ini
diberikan oleh distribusi $X_k$.
\par
Model lain dari jalan acak (terutama digunakan ketika daerah nilainya adalah ${\mathbf R}^1$) ialah sebuah
permainan yang melibatkan dua orang dan terdiri atas barisan langkah yang saling bebas dan
berdistribusi identik.  Jumlah $S_n$, misalnya, menyatakan skor orang pertama
setelah $n$ langkah, dengan asumsi bahwa skor orang kedua adalah $-S_n$.  Sebagai
contoh, dua orang dapat melambungkan koin, dengan kecocokan atau ketidakcocokan
masing-masing menyatakan $+1$ atau $-1$
bagi pemain pertama.  Atau, mungkin satu koin dilambungkan, dengan sisi kepala atau ekor
masing-masing menyatakan $+1$ atau $-1$ bagi pemain pertama.
\subsection*{Jalan Acak pada Garis Real}
Pertama-tama kita akan meninjau kasus taktrivial yang paling sederhana
dari jalan acak dalam ${\mathbf R}^1$, yaitu kasus ketika peubah acak $X_n$
memiliki fungsi distribusi yang sama, yang diberikan oleh
$$
f_X(x) = \left \{ \begin{array}{ll}
                        1/2, & \mbox{jika $x = \pm 1,$} \\
                          0, & \mbox{selain itu.} 
                  \end{array}
         \right.  
$$
Situasi ini bersesuaian dengan pelambungan sebuah koin seimbang, dengan $S_n$ menyatakan banyaknya
sisi kepala dikurangi banyaknya sisi ekor yang muncul dalam $n$ pelambungan pertama.  Perhatikan bahwa dalam
situasi ini, semua lintasan dengan panjang $n$ memiliki probabilitas yang sama, yaitu $2^{-n}$.
\par
Kadang-kadang akan membantu jika jalan acak digambarkan sebagai garis poligonal, atau lintasan, pada
bidang, dengan sumbu mendatar menyatakan waktu dan sumbu tegak menyatakan nilai
$S_n$.  Dengan diberikan barisan jumlah parsial $\{S_n\}$, mula-mula kita gambarkan titik-titik $(n, S_n)$, lalu
untuk setiap $k < n$, kita hubungkan $(k, S_k)$ dan $(k+1, S_{k+1})$ dengan ruas garis lurus.
\emx {Panjang} suatu lintasan hanyalah selisih nilai waktu antara titik awal
dan titik akhir lintasan.  Lihat Gambar \ref{fig 12.1}.  Gambar ini
dan proses yang diperagakannya identik dengan contoh pada Bab
1 tentang dua orang yang bermain kepala atau ekor.
\putfig{3.5truein}{PSfig1-3}{Jalan acak dengan panjang 40.}{fig 12.1}
\subsection*{Kembalinya dan Kembalinya yang Pertama}
Kita mengatakan bahwa suatu \emx {penyamaan}\index{penyamaan} telah terjadi, atau terdapat \emx {kembali
ke titik asal}\index{kembali ke titik asal} pada waktu $n$, jika $S_n = 0$.  Perhatikan bahwa hal ini hanya dapat
terjadi jika $n$ adalah bilangan bulat genap.  Untuk menghitung probabilitas penyamaan pada waktu $2m$,
kita hanya perlu menghitung banyaknya lintasan dengan panjang
$2m$ yang berawal dan berakhir di titik asal.  Banyaknya lintasan tersebut jelas
$${2m \choose m}\ .$$
Karena setiap lintasan memiliki probabilitas $2^{-2m}$, kita memperoleh teorema berikut.
\begin{theorem}\label{thm 12.1.1} 
Probabilitas kembali ke
titik asal pada waktu $2m$ diberikan oleh 
$$u_{2m} = {2m \choose m}2^{-2m}\ .$$
Probabilitas kembali ke titik asal pada waktu ganjil adalah 0.
\end{theorem}%\endthm
Suatu jalan acak dikatakan mengalami \emx {kembali pertama} ke titik asal\index{kembali pertama ke
titik asal} \index{kembali ke titik asal!pertama} pada waktu
$2m$ jika
$m > 0$, dan
$S_{2k} \ne 0$ untuk semua $k < m$.  Pada Gambar~\ref{fig 12.1}, kembali pertama terjadi pada waktu 2. 
Kita definisikan $f_{2m}$ sebagai probabilitas kejadian ini.  (Kita juga mendefinisikan $f_0 = 0$.)  Ungkapan
$f_{2m}2^{2m}$ dapat dipandang sebagai banyaknya lintasan dengan panjang
$2m$ antara titik $(0, 0)$ dan $(2m, 0)$ yang tidak menyentuh sumbu mendatar kecuali pada
kedua titik ujungnya.  Dengan menggunakan gagasan ini, teorema berikut mudah dibuktikan.
\begin{theorem}\label{thm 12.1.2} 
Untuk $n \ge 1$, probabilitas $\{u_{2k}\}$ dan $\{f_{2k}\}$ dihubungkan oleh persamaan
$$u_{2n} = f_0 u_{2n} + f_2 u_{2n-2} + \cdots + f_{2n}u_0\ .$$
\proof
Terdapat $u_{2n}2^{2n}$ lintasan dengan panjang $2n$ yang memiliki titik ujung $(0, 0)$ dan $(2n, 0)$.  Kumpulan
lintasan tersebut dapat dipartisi menjadi $n$ himpunan, bergantung pada waktu
kembali pertama ke titik asal.  Sebuah lintasan dalam kumpulan ini yang kembali untuk pertama kalinya ke titik asal
pada waktu $2k$ terdiri atas ruas awal dari $(0, 0)$ ke $(2k, 0)$, yang tidak memiliki titik dalam
pada sumbu mendatar, dan ruas akhir dari $(2k, 0)$ ke $(2n, 0)$, tanpa
batasan lebih lanjut pada ruas ini.  Jadi, banyaknya lintasan dalam kumpulan yang
kembali untuk pertama kalinya ke titik asal pada waktu $2k$ diberikan oleh
$$f_{2k}2^{2k}u_{2n-2k}2^{2n-2k} = f_{2k}u_{2n-2k}2^{2n}\ .$$
Jika kita menjumlahkan atas $k$, kita memperoleh persamaan
$$u_{2n}2^{2n} = f_0u_{2n} 2^{2n} + f_2u_{2n-2}2^{2n} + \cdots + f_{2n}u_0 2^{2n}\ .$$
Membagi kedua ruas persamaan ini dengan $2^{2n}$ menyelesaikan pembuktian.
\end{theorem}
Ungkapan pada ruas kanan teorema di atas seharusnya mengingatkan pembaca pada suatu jumlah
yang muncul dalam Definisi~\ref{defn 7.1} tentang konvolusi dua distribusi.  Konvolusi
dua barisan didefinisikan dengan cara serupa.  Teorema di atas menyatakan bahwa
barisan $\{u_{2n}\}$ adalah konvolusi dirinya sendiri dengan barisan
$\{f_{2n}\}$.  Jadi, jika setiap barisan ini kita nyatakan dengan fungsi pembangkit
biasa, kita dapat menggunakan hubungan di atas untuk menentukan nilai $f_{2n}$.
\begin{theorem}\label{thm 12.1.3}
Untuk $m \ge 1$, probabilitas kembali pertama ke titik asal pada waktu $2m$ diberikan oleh
$$f_{2m} = {{u_{2m}}\over{2m-1}} = {{2m \choose m}\over{(2m-1)2^{2m}}}\ .$$
\proof
Kita mulai dengan mendefinisikan fungsi pembangkit
$$U(x) = \sum_{m = 0}^\infty u_{2m}x^m$$
dan
$$F(x) = \sum_{m = 0}^\infty f_{2m}x^m\ .$$
Teorema~\ref{thm 12.1.2} menyatakan bahwa
\begin{equation} 
U(x) = 1 + U(x)F(x)\ .
\label{eq 12.1.1}  
\end{equation} 
(Adanya 1 pada ruas kanan disebabkan oleh fakta bahwa $u_0$ didefinisikan sama dengan 1,
tetapi Teorema~\ref{thm 12.1.2} hanya berlaku untuk $m \ge 1$.)  Perhatikan bahwa kedua fungsi pembangkit
tentu konvergen pada selang $(-1, 1)$ karena nilai mutlak semua koefisiennya paling besar 1.
Jadi, kita dapat menyelesaikan persamaan di atas untuk $F(x)$ dan memperoleh
$$F(x) = {{U(x) - 1}\over{U(x)}}\ .$$
Sekarang, jika kita dapat menemukan ungkapan bentuk tertutup untuk fungsi $U(x)$, kita juga akan memiliki
ungkapan bentuk tertutup untuk $F(x)$.  Dari Teorema~\ref{thm 12.1.1}, kita memperoleh
$$U(x) = \sum_{m = 0}^\infty {2m \choose m}2^{-2m}x^m\ .$$
Dalam Wilf,\footnote{H.~S.~Wilf, \emx {Generatingfunctionology,}\index{WILF, H. S.}
(Boston: Academic Press, 1990), hlm. 50.} kita mendapati bahwa
$${1\over{\sqrt {1 - 4x}}} = \sum_{m = 0}^\infty {2m \choose m} x^m\ .$$
Pembaca diminta membuktikan pernyataan ini dalam Latihan~\ref{exer 12.1.1}.  Jika $x$ kita ganti dengan $x/4$
dalam persamaan terakhir, kita melihat bahwa
$$U(x) = {1\over{\sqrt {1-x}}}\ .$$
Oleh karena itu, kita memperoleh
\begin{eqnarray*}
F(x) &=& {{U(x) - 1}\over{U(x)}}\\
     &=& {{(1-x)^{-1/2} - 1}\over{(1-x)^{-1/2}}}\\
     &=& 1 - (1-x)^{1/2}\ . 
\end{eqnarray*}

Meskipun nilai $f_{2m}$ dapat dihitung menggunakan Teorema Binomial, lebih
mudah untuk memperhatikan bahwa $F'(x) = U(x)/2$, sehingga koefisien $f_{2m}$ dapat ditemukan dengan
mengintegralkan deret untuk $U(x)$.  Untuk $m \ge 1$, kita memperoleh

\begin{eqnarray*}
f_{2m} &=& {u_{2m-2}\over{2m}}\\
       &=& {2m-2 \choose m-1}\over{m2^{2m-1}}\\
       &=& {2m \choose m}\over{(2m-1)2^{2m}}\\
       &=& {u_{2m}\over{2m-1}}\ ,
\end{eqnarray*}

karena
$${2m-2 \choose m-1} = {m\over{2(2m-1)}}{2m\choose m}\ .$$
Ini menyelesaikan pembuktian teorema.
\end{theorem}
\subsection*{Probabilitas Akhirnya Kembali}\index{kembali ke titik asal!probabilitas akhirnya}
Dalam proses jalan acak simetris di ${\mathbf R}^m$, berapakah probabilitas bahwa partikel
pada akhirnya kembali ke titik asal?  Pertama-tama kita telaah pertanyaan ini untuk kasus $m = 1$, lalu
kita tinjau kasus umumnya.  Hasil-hasil dalam dua contoh berikut berasal dari
P\'olya.\footnote{G.~P\'olya, ``\"Uber eine Aufgabe der Wahrscheinlichkeitsrechnung betreffend
die Irrfahrt im Strassennetz,'' Math. Ann., vol.~84 (1921), pp. 149-160.}\index{P\'OLYA, G.}
\begin{example}\label{exam 12.1.0.5}(Akhirnya Kembali dalam ${\mathbf R}^1$)
Gagasan tentang akhirnya kembali harus didekati dengan hati-hati karena
ruang sampelnya tampak berupa himpunan semua jalan dengan panjang takhingga, dan himpunan ini
tak terhitung.  Untuk menghindari kesulitan, kita definisikan $w_n$ sebagai probabilitas bahwa kembali
pertama telah terjadi selambat-lambatnya pada waktu $n$.  Jadi, $w_n$ berkaitan dengan ruang sampel semua jalan
dengan panjang $n$, yang merupakan himpunan berhingga.  Dengan memakai $w_n$, wajar untuk mendefinisikan
probabilitas bahwa partikel pada akhirnya kembali ke titik asal sebagai 
$$w_* = \lim_{n \rightarrow \infty} w_n\ .$$
Limit ini jelas ada dan paling besar satu karena barisan $\{w_n\}_{n = 1}^\infty$ adalah
barisan naik, dan semua sukunya paling besar satu.  
\par
Dengan menggunakan probabilitas $f_n$, kita melihat bahwa
$$w_{2n} = \sum_{i = 1}^n f_{2i}\ .$$
Jadi, 
$$w_* = \sum_{i = 1}^\infty f_{2i}\ .$$
Dalam pembuktian Teorema~\ref{thm 12.1.3}, fungsi pembangkit 
$$F(x) = \sum_{m = 0}^\infty f_{2m}x^m$$ 
diperkenalkan.  Di sana dinyatakan bahwa deret ini konvergen untuk $x \in (-1, 1)$.  Bahkan,
dapat ditunjukkan bahwa deret ini juga konvergen untuk $x = \pm 1$ dengan menggunakan Latihan~\ref{exer
12.1.4}, bersama dengan fakta bahwa
$$f_{2m} = {{u_{2m}}\over{2m-1}}\ .$$
(Fakta ini dibuktikan dalam pembuktian Teorema~\ref{thm 12.1.3}.)  Karena kita juga mengetahui bahwa 
$$F(x) = 1 - (1-x)^{1/2}\ ,$$
kita melihat bahwa 
$$w_* = F(1) = 1\ .$$
Jadi, partikel kembali ke titik asal dengan probabilitas satu.
\par
Bukti alternatif untuk fakta bahwa $w_* = 1$ dapat diperoleh dengan menggunakan
hasil-hasil dalam Latihan~\ref{exer 12.1.2}.
\end{example}
\par
\begin{example}\label{exam 12.1.0.6}(Akhirnya Kembali dalam ${\mathbf R}^m$)
Sekarang kita beralih ke kasus ketika jalan acak berlangsung dalam lebih dari satu dimensi.
Kita definisikan $f^{(m)}_{2n}$ sebagai probabilitas bahwa kembali pertama ke titik asal dalam ${\mathbf
R}^m$ terjadi pada waktu $2n$.  Besaran $u^{(m)}_{2n}$ didefinisikan dengan cara serupa.  Jadi,
$f^{(1)}_{2n}$ dan $u^{(1)}_{2n}$ masing-masing sama dengan $f_{2n}$ dan $u_{2n}$ yang telah didefinisikan sebelumnya.  Jika,
selain itu, kita mendefinisikan $u^{(m)}_0 = 1$ dan $f^{(m)}_0 = 0$, kita dapat meniru pembuktian
Teorema~\ref{thm 12.1.2} dan menunjukkan bahwa untuk semua $m \ge 1$,
\begin{equation}
u^{(m)}_{2n} = f^{(m)}_0 u^{(m)}_{2n} + f^{(m)}_2 u^{(m)}_{2n-2} + \cdots +
f^{(m)}_{2n}u^{(m)}_0\ .
\label{eq 12.1.1.5}
\end{equation}
\par
Kita melanjutkan generalisasi hasil sebelumnya dengan mendefinisikan
$$U^{(m)}(x) = \sum_{n = 0}^\infty u^{(m)}_{2n} x^n$$
dan
$$F^{(m)}(x) = \sum_{n = 0}^\infty f^{(m)}_{2n} x^n\ .$$
Kemudian, dengan menggunakan Persamaan~\ref{eq 12.1.1.5}, kita melihat bahwa
$$U^{(m)}(x) = 1 + U^{(m)}(x) F^{(m)}(x)\ ,$$
seperti sebelumnya.  Fungsi-fungsi ini selalu konvergen pada selang $(-1, 1)$ karena besar
semua koefisiennya paling besar satu.  Bahkan, karena 
$$w^{(m)}_* = \sum_{n = 0}^\infty f^{(m)}_{2n} \le 1$$
untuk semua $m$, deret untuk $F^{(m)}(x)$ juga konvergen pada $x = 1$, dan $F^{(m)}(x)$
kontinu dari kiri pada $x = 1$, yaitu 
$$\lim_{x \uparrow 1} F^{(m)}(x) = F^{(m)}(1)\ .$$
Jadi, kita memperoleh
\begin{equation}
w^{(m)}_* = \lim_{x \uparrow 1} F^{(m)}(x) = \lim_{x \uparrow 1} 
\frac{U^{(m)}(x) - 1}{U^{(m)}(x)}\ ,
\label{eq 12.1.1.6}
\end{equation}
sehingga untuk menentukan $w^{(m)}_*$, cukup dengan menentukan 
$$\lim_{x \uparrow 1} U^{(m)}(x)\ .$$
Kita nyatakan limit ini dengan $u^{(m)}$.
\par
Kita mengklaim bahwa 
$$u^{(m)} = \sum_{n = 0}^\infty u^{(m)}_{2n}\ .$$
(Klaim ini masuk akal; klaim ini menyatakan bahwa untuk mengetahui apa yang terjadi pada fungsi $U^{(m)}(x)$ di
$x = 1$, cukup tetapkan $x = 1$ dalam deret pangkat untuk $U^{(m)}(x)$.)  Untuk membuktikan klaim ini, perhatikan
bahwa koefisien $u^{(m)}_{2n}$ taknegatif, sehingga $U^{(m)}(x)$ naik secara monoton pada
selang $[0, 1)$.  Jadi, untuk setiap $K$, kita memperoleh
$$\sum_{n = 0}^K u^{(m)}_{2n} \le \lim_{x \uparrow 1} U^{(m)}(x) = u^{(m)} \le \sum_{n = 0}^\infty
u^{(m)}_{2n}\ .$$
Dengan mengambil $K \rightarrow \infty$, kita melihat bahwa
$$u^{(m)} = \sum_{2n}^\infty u^{(m)}_{2n}\ .$$
Ini membuktikan klaim tersebut.
\par
Dari Persamaan~\ref{eq 12.1.1.6}, kita melihat bahwa jika $u^{(m)} < \infty$, maka probabilitas
akhirnya kembali adalah
$$\frac {u^{(m)} - 1}{u^{(m)}}\ ,$$
sedangkan jika $u^{(m)} = \infty$, maka probabilitas akhirnya kembali adalah 1.
\par
Untuk menyelesaikan contoh ini, kita harus menaksir jumlah
$$\sum_{n = 0}^\infty u^{(m)}_{2n}\ .$$  
Dalam Latihan~\ref{exer 12.1.12}, pembaca diminta menunjukkan bahwa
$$u^{(2)}_{2n} = \frac 1 {4^{2n}} {{2n}\choose n}^2\ .$$
Dengan menggunakan Rumus Stirling, mudah ditunjukkan bahwa (lihat Latihan~\ref{exer 12.1.13})
$${{2n}\choose n} \sim \frac {2^{2n}}{\sqrt {\pi n}}\ ,$$
sehingga
$$u^{(2)}_{2n} \sim \frac 1{\pi n}\ .$$
Dari sini mudah disimpulkan bahwa
$$\sum_{n = 0}^\infty u^{(2)}_{2n}$$
divergen, sehingga $w^{(2)}_* = 1$; yaitu, dalam ${\mathbf R}^2$, probabilitas akhirnya kembali adalah
1.
\par
Ketika $m = 3$, Latihan~\ref{exer 12.1.12} menunjukkan bahwa
$$u^{(3)}_{2n} = \frac 1{2^{2n}}{{2n}\choose n} \sum_{j,k} 
\biggl(\frac 1{3^n}\frac{n!}{j!k!(n-j-k)!}\biggr)^2\ .$$
Misalkan $M$ menyatakan nilai terbesar dari 
$$\frac 1{3^n}\frac {n!}{j!k!(n - j - k)!}\ ,$$
di antara semua nilai taknegatif $j$ dan $k$ dengan $j + k \le n$.  Dengan menggunakan Rumus
Stirling, mudah ditunjukkan bahwa 
$$M \sim \frac cn\ ,$$
untuk suatu konstanta $c$.  Jadi, kita memperoleh
$$u^{(3)}_{2n} \le \frac 1{2^{2n}}{{2n}\choose n} \sum_{j,k} 
\biggl(\frac M{3^n}\frac{n!}{j!k!(n-j-k)!}\biggr)\ .$$
Dengan menggunakan Latihan~\ref{exer 12.1.14}, dapat ditunjukkan bahwa ungkapan pada ruas kanan paling besar
$$\frac {c'}{n^{3/2}}\ ,$$
dengan $c'$ suatu konstanta.  Jadi,
$$\sum_{n = 0}^\infty u^{(3)}_{2n}$$
konvergen, sehingga $w^{(3)}_*$ lebih kecil secara ketat dari satu.  Artinya, dalam ${\mathbf R}^3$,
probabilitas akhirnya kembali ke titik asal lebih kecil secara ketat dari satu (bahkan nilainya
sekitar .34).
\par
Hasil-hasil ini dapat dirangkum dengan menyatakan bahwa seseorang sebaiknya tidak mabuk dalam lebih dari dua
dimensi.
\end{example}

\subsection*{Nilai Harapan Banyaknya Penyamaan}
Sekarang kita berikan contoh lain penggunaan fungsi pembangkit untuk menemukan rumus umum bagi
suku-suku suatu barisan, ketika barisan tersebut dihubungkan dengan barisan lain melalui relasi rekurensi.
Latihan~\ref{exer 12.1.9} memberikan satu contoh lagi.
\begin{example}(Nilai Harapan Banyaknya Penyamaan)\label{exam
12.1.1}\index{penyamaan!nilai harapan banyaknya} 
Dalam contoh ini, kita akan menurunkan rumus untuk nilai harapan banyaknya penyamaan dalam jalan
acak dengan panjang $2m$.  Seperti dalam pembuktian Teorema~\ref{thm 12.1.3}, metode ini memiliki empat bagian utama. 
Pertama, ditemukan rekurensi yang menghubungkan suku ke-$m$ dalam barisan yang belum diketahui dengan
suku-suku sebelumnya dalam barisan yang sama serta suku-suku dalam barisan lain (yang diketahui).  Contoh
rekurensi semacam itu diberikan dalam Teorema~\ref{thm 12.1.2}.  Kedua, rekurensi tersebut digunakan untuk menurunkan
persamaan fungsional yang melibatkan fungsi pembangkit dari barisan yang belum diketahui dan satu atau lebih
barisan yang diketahui.  Persamaan~\ref{eq 12.1.1} merupakan contoh persamaan fungsional semacam itu.  Ketiga,
persamaan fungsional tersebut diselesaikan untuk fungsi pembangkit yang belum diketahui.  Terakhir, dengan menggunakan
cara seperti
Teorema Binomial, integrasi, atau diferensiasi, ditemukan rumus untuk koefisien ke-$m$
dari fungsi pembangkit yang belum diketahui.
\par
Kita mulai dengan mendefinisikan $g_{2m}$ sebagai banyaknya seluruh penyamaan di antara semua jalan acak
dengan panjang $2m$.  (Untuk setiap jalan acak, penyamaan pada waktu 0 kita abaikan.)  Kita
mendefinisikan $g_0 = 0$.  Karena banyaknya jalan dengan panjang $2m$ sama dengan $2^{2m}$, nilai harapan
banyaknya penyamaan di antara semua jalan acak tersebut adalah $g_{2m}/2^{2m}$.  Selanjutnya, kita mendefinisikan
fungsi pembangkit
$G(x)$:
$$G(x) = \sum_{k = 0}^\infty g_{2k}x^k\ .$$
Sekarang kita perlu menemukan rekurensi yang menghubungkan barisan $\{g_{2k}\}$ dengan salah satu atau kedua
barisan yang diketahui, $\{f_{2k}\}$ dan $\{u_{2k}\}$.  Kita anggap $m$ sebagai bilangan bulat positif yang tetap,
dan memandang himpunan semua lintasan dengan panjang $2m$ sebagai gabungan saling lepas
$$ E_2 \cup E_4 \cup \cdots \cup E_{2m} \cup H\ ,$$
dengan $E_{2k}$ adalah himpunan semua lintasan dengan panjang $2m$ yang mengalami penyamaan pertama pada waktu $2k$,
dan $H$ adalah himpunan semua lintasan dengan panjang $2m$ tanpa penyamaan.  Mudah ditunjukkan (lihat
Latihan~\ref{exer 12.1.3}) bahwa
$$|E_{2k}| = f_{2k} 2^{2m}\ .$$
Kita mengklaim bahwa banyaknya seluruh penyamaan pada semua lintasan yang termasuk dalam himpunan $E_{2k}$ sama
dengan 
\begin{equation}
|E_{2k}| + 2^{2k} f_{2k} g_{2m - 2k}\ .
\label{eq 12.1.2}  
\end{equation}
Setiap lintasan dalam $E_{2k}$ mengalami satu penyamaan pada waktu $2k$, sehingga jumlah seluruh
penyamaan tersebut tepat $|E_{2k}|$.  Ini adalah suku pertama dalam ungkapan Persamaan~\ref{eq 12.1.2}. 
Terdapat $2^{2k} f_{2k}$ ruas awal berbeda dengan panjang $2k$ di antara lintasan-lintasan dalam
$E_{2k}$.  Setiap ruas awal ini dapat diperpanjang menjadi lintasan dengan panjang $2m$ melalui
$2^{2m-2k}$ cara, dengan menyambungkan semua lintasan yang mungkin dengan panjang $2m - 2k$.  Berdasarkan
definisi, banyaknya penyamaan yang diperoleh dengan menyambungkan semua lintasan ini ke salah satu ruas awal adalah
$g_{2m - 2k}$.  Ini menghasilkan suku kedua dalam Persamaan~\ref{eq 12.1.2}.  Karena $k$
dapat berkisar dari 1 sampai $m$, kita memperoleh rekurensi
\begin{equation}  
g_{2m} = \sum_{k = 1}^m \Bigl(|E_{2k}| + 2^{2k}f_{2k}g_{2m - 2k}\Bigr)\ .
\label{eq 12.1.3}  
\end{equation}                                                                        
\par
Suku kedua dalam suku umum di atas seharusnya mengingatkan pembaca pada konvolusi.  Bahkan,
jika fungsi pembangkit $G(x)$ kita kalikan dengan fungsi pembangkit 
$$F(4x) = \sum_{k = 0}^\infty 2^{2k}f_{2k} x^k\ ,$$
koefisien $x^m$ sama dengan
$$\sum_{k = 0}^m 2^{2k}f_{2k}g_{2m-2k}\ .$$
Jadi, hasil kali $G(x)F(4x)$ merupakan bagian dari persamaan fungsional yang kita cari.  Suku
pertama dalam suku umum pada Persamaan~\ref{eq 12.1.3} menghasilkan jumlah
$$2^{2m}\sum_{k = 1}^m f_{2k}\ .$$
Dari Latihan~\ref{exer 12.1.2}, kita melihat bahwa jumlah ini tepat $(1 - u_{2m})2^{2m}$.    Jadi, kita perlu
membentuk fungsi pembangkit yang koefisien ke-$m$-nya adalah suku ini; fungsi pembangkit tersebut adalah
$$\sum_{m = 0}^\infty (1- u_{2m})2^{2m} x^m\ ,$$
atau
$$\sum_{m = 0}^\infty 2^{2m} x^m - \sum_{m = 0}^\infty u_{2m}2^{2m} x^m\ .$$
Jumlah pertama tepat $(1-4x)^{-1}$, dan jumlah kedua adalah $U(4x)$.  Jadi, persamaan
fungsional yang kita cari adalah
$$G(x) = F(4x)G(x) + {1\over{1-4x}} - U(4x)\ .$$
Jika rekurensi ini kita selesaikan untuk $G(x)$ lalu disederhanakan, kita memperoleh
\begin{equation}  
G(x) = {1\over{(1-4x)^{3/2}}} - {1\over{(1-4x)}}\ .
\label{eq 12.1.4}  
\end{equation}                                                                        
\par
Sekarang kita perlu menemukan rumus untuk koefisien $x^m$.  Suku pertama dalam 
Persamaan~\ref{eq 12.1.4} adalah $(1/2)U'(4x)$, sehingga koefisien $x^m$ dalam fungsi ini adalah
$$u_{2m+2} 2^{2m+1}(m+1)\ .$$
Suku kedua dalam Persamaan~\ref{eq 12.1.4} adalah jumlah deret geometri dengan rasio 
$4x$, sehingga koefisien $x^m$ adalah $2^{2m}$.  Jadi, kita memperoleh

\begin{eqnarray*}
g_{2m} &=& u_{2m+2}2^{2m+1}(m+1) - 2^{2m}\\
       &=& {1\over 2}{{2m+2} \choose {m+1}} (m+1)  - 2^{2m}\ .
\end{eqnarray*}

Ingat bahwa hasil bagi $g_{2m}/2^{2m}$ adalah nilai harapan banyaknya penyamaan pada semua
lintasan dengan panjang $2m$.  Dengan menggunakan Latihan~\ref{exer 12.1.4}, mudah ditunjukkan bahwa
$${{g_{2m}}\over{2^{2m}}} \sim \sqrt{2\over \pi}\sqrt {2m}\ .$$
Secara khusus, ini berarti bahwa rata-rata banyaknya penyamaan pada semua lintasan dengan panjang
$4m$ bukanlah dua kali rata-rata banyaknya penyamaan pada semua lintasan dengan panjang $2m$.  Agar
rata-rata banyaknya penyamaan menjadi dua kali lipat, panjang jalan acak harus dibuat empat kali lipat.
\end{example}
Menarik untuk diperhatikan bahwa jika kita mendefinisikan
$$M_n = \max_{0 \le k \le n} S_k\ ,$$
maka kita memperoleh
$$E(M_n) \sim \sqrt{2\over \pi}\sqrt n\ .$$
Artinya, nilai harapan banyaknya penyamaan dan nilai harapan dari nilai maksimum bagi jalan
acak dengan panjang $n$ sama secara asimtotik ketika $n \rightarrow \infty$.  (Bahkan, dapat
ditunjukkan bahwa kedua nilai harapan itu berbeda paling banyak $1/2$ untuk semua bilangan bulat positif $n$.  Lihat
Latihan~\ref{exer 12.1.9}.)

\exercises
\begin{LJSItem}
 
\i\label{exer 12.1.1} Dengan menggunakan Teorema Binomial, tunjukkan bahwa
$${1\over{\sqrt {1 - 4x}}} = \sum_{m = 0}^\infty {2m \choose m} x^m\ .$$
Berapakah interval konvergensi deret pangkat ini?

\i\label{exer 12.1.2}
\begin{enumerate}

\item Tunjukkan bahwa untuk $m \ge 1$, 
$$f_{2m} = u_{2m-2} - u_{2m}\ .$$

\item Dengan menggunakan bagian (a), carilah ekspresi bentuk tertutup untuk jumlah
$$f_2 + f_4 + \cdots + f_{2m}\ .$$

\item Dengan menggunakan bagian (b), tunjukkan bahwa
$$\sum_{m = 1}^\infty f_{2m} = 1\ .$$
(Pernyataan ini juga dapat diperoleh dari kenyataan bahwa 
$$F(x) = 1 - (1-x)^{1/2}\ .)$$

\item Dengan menggunakan bagian (a) dan (b), tunjukkan bahwa peluang tidak
terjadi penyamaan dalam $2m$ hasil pertama sama dengan peluang terjadinya
penyamaan pada waktu $2m$.
\end{enumerate}

\i\label{exer 12.1.3} Dengan menggunakan notasi Contoh~\ref{exam 12.1.1},
tunjukkan bahwa
$$|E_{2k}| = f_{2k} 2^{2m}\ .$$

\i\label{exer 12.1.4} Dengan menggunakan Rumus Stirling, tunjukkan bahwa 
$$u_{2m} \sim {1\over{\sqrt {\pi m}}}\ .$$

\i\label{exer 12.1.5} Suatu \emx {pergantian keunggulan}\index{pergantian
keunggulan} dalam perjalanan acak terjadi pada waktu $2k$ jika $S_{2k-1}$ dan
$S_{2k+1}$ mempunyai tanda berlawanan.  
\begin{enumerate}
\item Berikan argumen ketat yang membuktikan bahwa di antara semua perjalanan
sepanjang $2m$ yang mengalami penyamaan pada waktu $2k$, tepat setengahnya
mengalami pergantian keunggulan pada waktu $2k$.

\item Simpulkan bahwa jumlah seluruh pergantian keunggulan di antara semua
perjalanan sepanjang $2m$ sama dengan
$${1\over 2}(g_{2m} - u_{2m})\ .$$

\item Carilah ekspresi asimtotik untuk banyaknya rata-rata pergantian keunggulan
dalam perjalanan acak sepanjang $2m$. 
\end{enumerate}

\i\label{exer 12.1.6} 
\begin{enumerate}
\item 
Tunjukkan bahwa peluang perjalanan acak sepanjang $2m$ kembali ke titik asal
untuk terakhir kalinya\index{kembali terakhir ke titik asal}\index{kembali ke
titik asal!terakhir} pada waktu $2k$, dengan $0 \le k \le m$, sama dengan
$$
{{{2k}\choose k}{{2m-2k}\choose {m-k}}\over{2^{2m}}} = u_{2k}u_{2m - 2k}\ .
$$  
(Kasus $k = 0$ terdiri atas semua lintasan yang tidak kembali ke titik asal pada
waktu positif mana pun.) \emx {Petunjuk}: Lintasan yang terakhir kembali ke
titik asal pada waktu $2k$ terdiri atas dua lintasan yang disambungkan: satu
lintasan sepanjang $2k$ yang berawal dan berakhir di titik asal, dan lintasan
lain sepanjang $2m - 2k$ yang berawal di titik asal tetapi tidak pernah kembali
ke titik asal. Kedua jenis lintasan dapat dihitung menggunakan besaran yang
muncul dalam bagian ini.  
\item
Dengan menggunakan bagian (a), tunjukkan bahwa jika $m$ ganjil, peluang bahwa
perjalanan sepanjang $2m$ tidak mengalami penyamaan dalam $m$ hasil terakhir
sama dengan $1/2$, berapa pun nilai $m$. \emx {Petunjuk}: Jawaban bagian a)
simetris dalam $k$ dan $m-k$.
\end{enumerate}

\i\label{exer 12.1.7} Tunjukkan bahwa peluang tidak terjadi penyamaan dalam
perjalanan sepanjang $2m$ sama dengan $u_{2m}$.

\istar\label{exer 12.1.8} Tunjukkan bahwa
$$P(S_1 \ge 0,\ S_2 \ge 0,\ \ldots,\ S_{2m} \ge 0) = u_{2m}\ .$$
\emx {Petunjuk}: Pertama, jelaskan mengapa
\begin{eqnarray*}
&&P(S_1 > 0,\ S_2 > 0,\ \ldots,\ S_{2m} > 0) \\
&& \;\;\;\;\;\;\;\;\;\;\;\;\; = {1\over 2}P(S_1 \ne 0,\ S_2 \ne 0,\ \ldots,\ S_{2m} \ne 0)  \ .
\end{eqnarray*}
Kemudian gunakan Latihan~\ref{exer 12.1.7}, bersama pengamatan bahwa jika tidak
terjadi penyamaan dalam $2m$ hasil pertama, maka lintasan melewati titik $(1,1)$
dan tetap berada pada atau di atas garis horizontal $x = 1$.

\istar\label{exer 12.1.9} Dalam Feller,\footnote{W.~Feller, \emx {Introduction to Probability
Theory and its Applications,} vol.~I, 3rd~ed. (New York: John Wiley \& Sons, 1968).}
terdapat teorema berikut: misalkan $M_n$ adalah peubah acak yang memberikan
nilai maksimum $S_k$, untuk $1 \le k \le n$. Definisikan
$$p_{n, r} = {n\choose {{n+r}\over{2}}}2^{-n}\ .$$
Jika $r \ge 0$, maka
$$P(M_n = r) = \left \{ \begin{array}{ll}
                    p_{n, r}\,,&\mbox{jika $r \equiv n\,    (\mbox{mod}\ 2)$}, \\   
                  p_{n, r+1}\,,&\mbox{jika $r \not\equiv n\,(\mbox{mod}\ 2)$}.
                      \end{array}
             \right.   
$$

\begin{enumerate}

\item Dengan menggunakan teorema ini, tunjukkan bahwa
$$E(M_{2m}) = {1\over{2^{2m}}}\sum_{k = 1}^m (4k-1){2m \choose m+k}\ ,$$
dan jika $n = 2m+1$, maka
$$E(M_{2m+1}) = {1\over {2^{2m+1}}} \sum_{k = 0}^m (4k+1){2m+1\choose m+k+1}\ .$$

\item
Untuk $m \ge 1$, definisikan 
$$r_m = \sum_{k = 1}^m k {2m\choose m+k}$$
dan
$$s_m = \sum_{k = 1}^m k {2m+1\choose m+k+1}\ .$$
Dengan menggunakan identitas
$${n\choose k} = {n-1\choose k-1} + {n-1\choose k}\ ,$$
tunjukkan bahwa
$$s_m = 2r_m - {1\over 2}\biggl(2^{2m} - {2m \choose m}\biggr)$$
dan
$$r_m = 2s_{m-1} + {1\over 2}2^{2m-1}\ ,$$
jika $m \ge 2$.

\item
Definisikan fungsi-fungsi pembangkit
$$R(x) = \sum_{k = 1}^\infty r_k x^k$$
dan
$$S(x) = \sum_{k = 1}^\infty s_k x^k\ .$$
Tunjukkan bahwa
$$S(x) = 2 R(x) - {1\over 2}\biggl({1\over{1- 4x}}\biggr) + {1\over 2}\biggl(\sqrt{1-4x}\biggr)$$
dan
$$R(x) = 2xS(x) + x\biggl({1\over{1-4x}}\biggr)\ .$$

\item
Tunjukkan bahwa
$$R(x) = {x\over{(1-4x)^{3/2}}}\ ,$$
dan
$$S(x) = {1\over 2}\biggl({1\over{(1- 4x)^{3/2}}}\biggr) - 
{1\over 2}\biggl({1\over{1- 4x}}\biggr)\ .$$

\item
Tunjukkan bahwa
$$r_m = m{2m-1\choose m-1}\ ,$$
dan
$$s_m = {1\over 2}(m+1){2m+1\choose m} - {1\over 2}(2^{2m})\ .$$

\item
Tunjukkan bahwa
$$E(M_{2m}) = {m\over{2^{2m-1}}}{2m\choose m} + {1\over{2^{2m+1}}}{2m\choose m} - {1\over 2}\ ,$$
dan
$$E(M_{2m+1}) = {{m+1}\over{2^{2m+1}}}{2m+2\choose m+1} - {1\over 2}\ .$$
Pembaca hendaknya membandingkan rumus-rumus ini dengan ekspresi untuk
\linebreak $g_{2m}/2^{(2m)}$ dalam Contoh~\ref{exam 12.1.1}.
\end{enumerate}

\istar\label{exer 12.1.10}
(dari K. Levasseur\footnote{K. Levasseur, ``How to Beat Your Kids at Their
Own Game," \emx {Mathematics Magazine} vol.~61, no. 5 (December, 1988),
pp.~301-305.})\index{LEVASSEUR, K.} Salah satu orang tua dan anaknya memainkan
permainan berikut. Satu tumpukan berisi $2n$ kartu, $n$ merah dan $n$ hitam,
dikocok. Kartu dibuka satu per satu. Sebelum setiap kartu dibuka, orang tua dan
anak menebak apakah kartu itu merah atau hitam. Siapa yang membuat lebih banyak
tebakan benar memenangkan permainan. Anak diasumsikan menebak masing-masing
warna dengan peluang yang sama, sehingga rata-rata skornya adalah $n$. Orang
tua mencatat berapa banyak kartu setiap warna yang telah dibuka. Jika kartu
hitam yang tersisa dalam tumpukan lebih banyak daripada kartu merah, misalnya,
orang tua menebak hitam; sedangkan jika jumlah kartu yang tersisa untuk setiap
warna sama, orang tua menebak masing-masing warna dengan peluang 1/2. Berapakah
banyaknya tebakan benar yang diharapkan dari orang tua? \emx {Petunjuk}: Setiap
dari ${{2n}\choose n}$ urutan kartu merah dan hitam yang mungkin bersesuaian
dengan perjalanan acak sepanjang $2n$ yang kembali ke titik asal pada waktu
$2n$. Tunjukkan bahwa di antara setiap pasangan penyamaan berturut-turut, orang
tua akan benar tepat satu kali lebih banyak daripada salah. Jelaskan mengapa
ini berarti banyaknya rata-rata tebakan benar oleh orang tua melebihi $n$
sebesar tepat setengah banyaknya rata-rata penyamaan. Sekarang definisikan
peubah acak $X_i$ bernilai 1 jika terjadi penyamaan pada waktu $2i$, dan 0 jika
tidak. Maka, di antara semua lintasan yang relevan, kita mempunyai
$$E(X_i) = P(X_i = 1) = \frac{{{2n-2i}\choose{n-i}}{{2i}\choose{i}}}{{{2n}\choose{n}}}\ .$$
Jadi, banyaknya penyamaan yang diharapkan sama dengan
$$E\biggl(\sum_{i = 1}^n X_i\biggr) = \frac 1{{{2n}\choose{n}}}\sum_{i = 1}^n 
{{2n-2i}\choose{n-i}}{{2i}\choose{i}}\ .$$
Sekarang fungsi pembangkit dapat digunakan untuk menemukan nilai jumlah tersebut.
\par
Perlu diperhatikan bahwa dalam permainan seperti ini, pertanyaan yang lebih
menarik daripada pertanyaan di atas adalah: berapakah peluang orang tua
memenangkan permainan? Untuk permainan ini, pertanyaan tersebut dijawab oleh
D. Zagier.\footnote{D. Zagier, ``How Often Should You Beat Your Kids?"
\emx{Mathematics Magazine} vol.~63, no.\ 2 (April 1990), pp. 89-92.}
\index{ZAGIER, D.} Ia menunjukkan bahwa peluang menang secara asimtotik (untuk
$n$ besar) sama dengan besaran
$$\frac 12 + \frac 1{2\sqrt 2}\ .$$

\item\label{exer 12.1.11}
Buktikan bahwa
$$u^{(2)}_{2n} = \frac 1{4^{2n}} \sum_{k = 0}^n \frac {(2n)!}{k!k!(n-k)!(n-k)!}\ ,$$
dan
$$u^{(3)}_{2n} = \frac 1{6^{2n}} \sum_{j,k} \frac {(2n)!}{j!j!k!k!(n-j-k)!(n-j-k)!}\ ,$$
di mana jumlah terakhir mencakup semua $j$ dan $k$ tak negatif dengan
$j+k \le n$. Tunjukkan pula bahwa ekspresi terakhir ini dapat ditulis ulang sebagai
$$\frac 1{2^{2n}}{{2n}\choose n} \sum_{j,k} 
\biggl(\frac 1{3^n}\frac{n!}{j!k!(n-j-k)!}\biggr)^2\ .$$

\item\label{exer 12.1.12}
Buktikan bahwa jika $n \ge 0$, maka
$$\sum_{k = 0}^n {n \choose k}^2 = {{2n} \choose n}\ .$$
\emx {Petunjuk}: Tuliskan jumlah itu sebagai 
$$\sum_{k = 0}^n {n \choose k}{n \choose {n-k}}$$
dan jelaskan mengapa ini merupakan koefisien dalam hasil kali
$$(1 + x)^n (1 + x)^n\ .$$
Gunakan hasil ini bersama Latihan~\ref{exer 12.1.11} untuk menunjukkan bahwa
$$u^{(2)}_{2n} = \frac 1{4^{2n}}{{2n}\choose n}\sum_{k = 0}^n {n \choose k}^2 = 
\frac 1 {4^{2n}} {{2n}\choose n}^2\ .$$

\item\label{exer 12.1.13}
Dengan menggunakan Rumus Stirling, buktikan bahwa
$${{2n}\choose n} \sim \frac {2^{2n}}{\sqrt {\pi n}}\ .$$

\item\label{exer 12.1.14}
Buktikan bahwa
$$\sum_{j,k} 
\biggl(\frac 1{3^n}\frac{n!}{j!k!(n-j-k)!}\biggr) = 1\ ,$$
di mana jumlah mencakup semua $j$ dan $k$ tak negatif sedemikian rupa sehingga
$j + k \le n$. \emx {Petunjuk}: Hitunglah banyaknya cara menempatkan $n$ bola
berlabel ke dalam 3 guci berlabel.

\item\label{exer 12.1.15}
Dengan menggunakan hasil yang dibuktikan untuk perjalanan acak dalam
${\mathbf R}^3$ pada Contoh~\ref{exam 12.1.0.6}, jelaskan mengapa peluang pada
akhirnya kembali dalam ${\mathbf R}^n$ benar-benar kurang dari satu untuk
semua $n \ge 3$. \emx {Petunjuk}: Perhatikan perjalanan acak dalam
${\mathbf R}^n$ dan abaikan semua kecuali tiga koordinat pertama posisi partikel.

\end{LJSItem}

\section{Kebangkrutan Penjudi\label{sec 12.2}}\index{kebangkrutan penjudi}
\par
Pada bagian sebelumnya, kita mempelajari jenis perjalanan acak simetris paling
sederhana dalam ${\mathbf R}^1$. Dalam bagian ini, kita menghapus asumsi bahwa
perjalanan acak itu simetris. Sebagai gantinya, kita mengandaikan bahwa $p$ dan
$q$ adalah bilangan real tak negatif dengan $p+q = 1$, dan bahwa fungsi
distribusi bersama lompatan-lompatan perjalanan acak adalah
$$f_X(x) = \left \{ \begin{array}{ll}
               p, & \mbox{jika $x = 1$},\\ 
               q, & \mbox{jika $x = -1$}.
                      \end{array}
             \right.   
$$ 
Perjalanan acak dapat dibayangkan sebagai representasi urutan lemparan koin
berbobot, dengan sisi kepala muncul dengan peluang $p$ dan sisi ekor muncul
dengan peluang $q$. Perumusan alternatif untuk situasi ini adalah seorang
penjudi yang memainkan serangkaian permainan melawan lawan (kadang-kadang
dianggap sebagai orang lain, kadang-kadang disebut ``bandar"), dengan peluang
$p$ bagi penjudi untuk menang dalam setiap permainan.

\subsection*{Masalah Kebangkrutan Penjudi}

Perumusan jenis perjalanan acak di atas menghasilkan masalah yang dikenal
sebagai masalah Kebangkrutan Penjudi.\index{kebangkrutan penjudi} Masalah ini
diperkenalkan dalam Latihan~\ref{exer 11.2.22}, tetapi kita akan menjelaskannya
kembali. Seorang penjudi memulai dengan ``modal" sebesar~$s$. Penjudi bermain
sampai modalnya mencapai nilai $M$ atau nilai 0. Dalam bahasa rantai Markov,
kedua nilai ini bersesuaian dengan keadaan penyerap. Kita tertarik mempelajari
peluang terjadinya masing-masing dari kedua hasil ini.  
\par
Kita juga dapat mengandaikan bahwa penjudi bermain melawan lawan yang ``kaya
tanpa batas.'' Dalam kasus ini, kita mengatakan bahwa hanya ada satu keadaan
penyerap, yaitu ketika modal penjudi bernilai 0. Dengan asumsi ini, kita dapat
menanyakan peluang bahwa penjudi pada akhirnya bangkrut. 
\par
Kita mulai dengan mendefinisikan $q_k$ sebagai peluang bahwa modal penjudi
mencapai 0, yakni ia bangkrut, sebelum mencapai $M$, dengan syarat modal awalnya
adalah $k$. Perhatikan bahwa $q_0 = 1$ dan $q_M = 0$. Relasi fundamental di
antara $q_k$ adalah sebagai berikut:
$$q_k = pq_{k+1} + qq_{k-1}\ ,$$
di mana $1 \le k \le M-1$. Relasi ini berlaku karena jika modalnya sama dengan
$k$ dan ia memainkan satu permainan, modalnya menjadi $k+1$ dengan peluang $p$
dan $k-1$ dengan peluang $q$. Dalam kasus pertama, peluang kebangkrutan pada
akhirnya adalah $q_{k+1}$ dan dalam kasus kedua adalah $q_{k-1}$. Karena
$p + q = 1$, persamaan di atas dapat kita tulis sebagai
$$p(q_{k+1} - q_k) = q(q_k - q_{k-1})\ ,$$
atau
$$q_{k+1} - q_k = {q\over p}(q_k - q_{k-1})\ .$$
Dari persamaan ini, mudah dilihat bahwa
\begin{equation}
q_{k+1} - q_k = \biggl({q\over p}\biggr)^k(q_1 - q_0)\ .
\label{eq 12.2.2}
\end{equation}
Sekarang kita menggunakan jumlah teleskopik untuk memperoleh persamaan yang
satu-satunya peubah tak dikenalnya adalah $q_1$:
\begin{eqnarray*}
-1 &=& q_M - q_0 \\
   &=& \sum_{k = 0}^{M-1} (q_{k+1} - q_k)\ ,
\end{eqnarray*}
sehingga
\begin{eqnarray*}
-1 &=& \sum_{k = 0}^{M-1} \biggl({q\over p}\biggr)^k(q_1 - q_0)\\
   &=& (q_1 - q_0) \sum_{k = 0}^{M-1} \biggl({q\over p}\biggr)^k\ .
\end{eqnarray*}
Jika $p \ne q$, maka ekspresi di atas sama dengan
$$(q_1 - q_0) {{(q/p)^M - 1}\over{(q/p) - 1}}\ ,$$
sedangkan jika $p = q = 1/2$, kita memperoleh persamaan
$$-1 = (q_1 - q_0) M\ .$$
Untuk sementara kita andaikan bahwa $p \ne q$. Maka kita mempunyai
$$q_1 - q_0 = -{{(q/p) - 1}\over{(q/p)^M - 1}}\ .$$
Sekarang, untuk setiap $z$ dengan $1 \le z \le M$, kita mempunyai
\begin{eqnarray*}
q_z - q_0 &=& \sum_{k = 0}^{z-1} (q_{k+1} - q_k)\\
          &=& (q_1 - q_0)\sum_{k = 0}^{z-1} \biggl({q\over p}\biggr)^k\\
          &=& -(q_1 - q_0){{(q/p)^z - 1}\over{(q/p) - 1}}\\
          &=& -{{(q/p)^z - 1}\over{(q/p)^M - 1}}\ .
\end{eqnarray*}
Oleh karena itu,
\begin{eqnarray*}
q_z &=& 1 - {{(q/p)^z - 1}\over{(q/p)^M - 1}}\\
    &=& {{(q/p)^M - (q/p)^z}\over{(q/p)^M - 1}}\ .
\label{eq 12.2.3}
\end{eqnarray*}
Akhirnya, jika $p = q = 1/2$, mudah ditunjukkan bahwa (lihat
Latihan~\ref{exer 12.2.11})
$$q_z = {{M - z}\over M}\ .$$
Perhatikan bahwa kedua rumus ini berlaku jika $z = 0$.
\par
Untuk $0 \le z \le M$, kita definisikan besaran $p_z$ sebagai peluang bahwa
modal penjudi mencapai $M$ tanpa pernah mencapai 0. Karena permainan mungkin
berlanjut tanpa batas, tidak jelas bahwa $p_z + q_z = 1$ untuk setiap $z$.
Namun, metode yang sama seperti di atas dapat digunakan untuk menunjukkan bahwa
jika $p \ne q$, maka
$$q_z = {{(q/p)^z - 1}\over{(q/p)^M - 1}}\ ,$$
dan jika $p = q = 1/2$, maka
$$q_z = {z\over M}\ .$$
Jadi, untuk setiap $z$, berlaku $p_z + q_z = 1$, sehingga permainan berakhir
dengan peluang 1.

\subsection*{Lawan yang Kaya Tanpa Batas}

Sekarang kita beralih ke masalah mencari peluang kebangkrutan pada akhirnya
jika penjudi bermain melawan lawan yang kaya tanpa batas. Peluang ini dapat
diperoleh dengan membiarkan $M$ menuju $\infty$ dalam ekspresi untuk $q_z$ yang
dihitung di atas. Jika $q < p$, ekspresi tersebut mendekati $(q/p)^z$, dan jika
$q > p$, ekspresi itu mendekati 1. Dalam kasus $p = q = 1/2$, ingat bahwa
$q_z = 1 - z/M$. Jadi, jika $M \rightarrow \infty$, kita melihat bahwa peluang
kebangkrutan pada akhirnya menuju 1.  

\subsection*{Catatan Historis}

Pada 1711, De Moivre\index{de MOIVRE, A.}, dalam bukunya \emx{De Mensura Sortis},
memberikan penurunan yang cerdik untuk peluang kebangkrutan. Ringkasan argumennya
berikut ini didasarkan pada uraian David.\footnote{F.~N. David,
\emx {Games, Gods and Gambling} (London: Griffin, 1962). Ringkasan editorial;
ungkapan sumber tidak diterjemahkan atau direproduksi.}\index{DAVID, F. N.}
Notasi yang digunakan adalah sebagai berikut: kita bayangkan terdapat dua
pemain, A dan B, dan peluang mereka memenangkan suatu permainan masing-masing
adalah $p$ dan $q$. Para pemain memulai dengan masing-masing $a$ dan $b$ keping.
Berikan nilai nominal $q/p$ pada keping terbawah A, $(q/p)^2$ pada keping di
atasnya, dan seterusnya hingga $(q/p)^a$ pada keping teratas A. Pada tumpukan B,
nilai keping teratas ialah $(q/p)^{a+1}$ dan nilai berikutnya berturut-turut
hingga $(q/p)^{a+b}$ pada keping terbawah. Sesudah setiap permainan, keping
teratas pihak yang kalah dipindahkan ke atas tumpukan pemenang; keping teratas
selalu menjadi taruhan permainan berikutnya. Dalam nilai nominal, taruhan B
selalu $q/p$ kali taruhan A, sehingga harapan nominal setiap pemain pada setiap
permainan adalah nol. Karena sifat ini berlaku sepanjang permainan, peluang A
memenangkan seluruh keping B dikalikan keuntungan nominal A harus sama dengan
peluang B dikalikan keuntungan nominal B. Maka
\begin{equation}
P_a\biggl(\Bigl({q\over p}\Bigr)^{a+1} + \cdots +
\Bigl({q\over p}\Bigr)^{a+b}\biggr) =  
P_b\biggl(\Bigl({q\over p}\Bigr) + \cdots +
\Bigl({q\over p}\Bigr)^a\biggr)\ .
\label{eq 12.2.1}
\end{equation}
\par
Dengan menggunakan persamaan ini bersama kenyataan bahwa
$$P_a + P_b = 1\ ,$$
dapat dengan mudah ditunjukkan bahwa
$$P_a = {{(q/p)^a - 1}\over{(q/p)^{a+b} - 1}}\ ,$$
jika $p \ne q$, dan
$$P_a = {a\over{a+b}}\ ,$$
jika $p = q = 1/2$.
\par
Dalam teori peluang modern, de Moivre mengubah nilai keping-keping untuk
mengubah permainan tak adil menjadi permainan adil, yang disebut martingal.
Dengan nilai baru itu, kekayaan harapan pemain A (yakni, jumlah nilai nominal
keping-kepingnya) setelah setiap permainan sama dengan kekayaannya sebelum
permainan (dan demikian pula untuk pemain B). (Untuk argumen martingal yang lebih
sederhana, lihat Latihan~\ref{exer 12.2.10}.) De Moivre kemudian menggunakan
kenyataan bahwa ketika permainan berakhir, permainan tetap adil, sehingga
Persamaan~\ref{eq 12.2.1} harus benar. Kenyataan ini memerlukan bukti dan
merupakan salah satu teorema sentral dalam teori martingal.  

\exercises
\begin{LJSItem}
 
\i\label{exer 12.2.2} Dalam masalah kebangkrutan penjudi, andaikan modal awal
penjudi adalah 1 dolar, dan andaikan peluangnya untuk menang dalam setiap
permainan adalah $p$. Misalkan $T$ adalah banyaknya permainan hingga modal 0
tercapai (penjudi bangkrut). Tunjukkan bahwa fungsi pembangkit untuk~$T$ adalah
$$
h(z) = \frac{1 - \sqrt{1 - 4pqz^2}}{2pz}\ ,
$$
dan bahwa
$$
h(1) = \left \{ \begin{array}{ll}
              q/p, & \mbox{jika $q \leq p$}, \\
                1, & \mbox{jika $q \geq p,$} 
         \end{array}
         \right.    
$$
dan
$$
h'(1) = \left \{ \begin{array}{ll}
               1/(q - p), & \mbox{jika $q > p$}, \\
                  \infty, & \mbox{jika $q = p.$}
         \end{array}
         \right. 
$$
Tafsirkan hasil Anda dalam kaitannya dengan waktu $T$ untuk mencapai 0. (Lihat
juga Contoh~\ref{exam 10.1.7}.)

\i\label{exer 12.2.3} Tunjukkan bahwa ekspansi deret Taylor untuk $\sqrt{1 - x}$
adalah
$$
 \sqrt{1 - x} = \sum_{n = 0}^\infty  {{1/2} \choose n} x^n\ ,
$$
di mana koefisien binomial ${1/2} \choose n$ adalah
$$
{{1/2} \choose n} = \frac{(1/2)(1/2 - 1) \cdots (1/2 - n + 1)}{n!}\ .  
$$
Dengan menggunakan hasil ini dan hasil Latihan~\ref{exer 12.2.2}, tunjukkan
bahwa peluang penjudi bangkrut pada langkah ke-$n$ adalah
$$
p_T(n) = \left \{ \begin{array}{ll}
 \frac{(-1)^{k - 1}}{2p} {{1/2} \choose k} (4pq)^k, & \mbox{jika $n = 2k - 1$,} \\
                                                 0, & \mbox{jika $n = 2k$.}
                  \end{array}
         \right.   
$$

\i\label{exer 12.2.4} Untuk masalah kebangkrutan penjudi, andaikan penjudi
memulai dengan $k$ dolar. Misalkan $T_k$ adalah waktu untuk mencapai 0 untuk
pertama kalinya.

\begin{enumerate}
\item Tunjukkan bahwa fungsi pembangkit $h_k(t)$ untuk~$T_k$ adalah pangkat
ke-$k$ dari fungsi pembangkit untuk waktu $T$ menuju kebangkrutan ketika
dimulai dari~1. \emx{Petunjuk}: Misalkan $T_k = U_1 + U_2 +\cdots+ U_k$, di
mana $U_j$ adalah waktu yang diperlukan perjalanan yang dimulai dari~$j$
untuk mencapai $j - 1$ untuk pertama kalinya.

\item Tentukan $h_k(1)$ dan $h_k'(1)$, lalu tafsirkan hasil Anda.
\end{enumerate} 

\i\label{exer 12.2.5} (Tiga soal berikut berasal dari Feller.\footnote{W.~Feller,
op.\ cit., pg. 367.}) Seperti dalam teks, andaikan $M$ adalah bilangan bulat
positif yang tetap. 

\begin{enumerate}
\item Tunjukkan bahwa jika seorang penjudi memulai dengan modal 0 (dan
diizinkan memiliki jumlah uang negatif), maka peluang modalnya mencapai nilai
$M$ sebelum kembali ke 0 sama dengan $p(1 - q_1)$.
\item Tunjukkan bahwa jika penjudi memulai dengan modal $M$, maka peluang
modalnya mencapai 0 sebelum kembali ke $M$ sama dengan $qq_{M-1}$.
\end{enumerate}

\i\label{exer 12.2.6} Misalkan seorang penjudi memulai dengan modal 0 dolar.
\begin{enumerate} 
\item Tunjukkan bahwa peluang modalnya tidak pernah mencapai $M$ sebelum
kembali ke 0 sama dengan $1 - p(1 - q_1)$.

\item Tunjukkan bahwa peluang modalnya mencapai nilai $M$ tepat $k$ kali
sebelum kembali ke 0 sama dengan
$p(1-q_1)(1 - qq_{M-1})^{k-1}(qq_{M-1})$. \emx {Petunjuk}: Gunakan
Latihan~\ref{exer 12.2.5}.
\end{enumerate} 

\i\label{exer 12.2.7} Dalam teks telah ditunjukkan bahwa jika $q < p$, terdapat
peluang positif bahwa penjudi yang memulai dengan modal 0 dolar tidak akan
pernah kembali ke titik asal. Karena itu, sekarang kita andaikan $q \ge p$.
Dengan menggunakan Latihan~\ref{exer 12.2.6}, tunjukkan bahwa jika seorang
penjudi memulai dengan modal 0 dolar, maka nilai harapan banyaknya kali
modalnya sama dengan $M$ sebelum kembali ke 0 adalah $(p/q)^M$ jika $q > p$,
dan 1 jika $q = p$. Feller memberikan tafsiran berikut: lempar koin seimbang
hingga banyak kumulatif sisi kepala dan sisi ekor pertama kali kembali sama, lalu
bayarkan satu sen setiap kali keunggulan kumulatif sisi kepala tepat $m$. Nilai
harapan pembayaran---dan karenanya biaya masuk yang adil---ialah satu sen untuk
setiap $m$.\footnote{W.~Feller, op.\ cit., p.~367. Ringkasan editorial; ungkapan
sumber tidak diterjemahkan atau direproduksi.}

\i\label{exer 12.2.8} Dalam permainan pada Latihan~\ref{exer 12.2.7}, ambil
$p = q = 1/2$ dan $M = 10$. Berapakah peluang bahwa modal penjudi sama dengan
$M$ setidaknya 20 kali sebelum kembali ke 0?

\i\label{exer 12.2.9} Tulis program komputer yang menyimulasikan permainan pada
Latihan~\ref{exer 12.2.7} untuk kasus $p = q = 1/2$ dan $M = 10$.  

\i\label{exer 12.2.10} Dalam uraian permainan oleh de Moivre, kita dapat
mengubah definisi kekayaan pemain A sedemikian rupa sehingga permainan tetap
merupakan martingal (dan perhitungannya menjadi lebih sederhana). Caranya ialah
memberikan nilai nominal kepada keping-keping seperti yang dilakukan de Moivre,
tetapi kekayaan setiap pemain saat ini didefinisikan hanya sebagai nilai keping
yang akan dipertaruhkan pada permainan berikutnya. Jadi, jika pemain A memiliki
$a$ keping, kekayaannya saat ini adalah $(q/p)^a$ (kita menetapkan bahwa ini
berlaku bahkan jika $a = 0$). Tunjukkan bahwa menurut definisi ini, jika
$p \ne q$, nilai harapan kekayaan pemain A setelah satu permainan sama dengan
kekayaannya sebelum permainan. Kemudian, seperti de Moivre, tuliskan persamaan
yang menyatakan bahwa nilai harapan kekayaan akhir pemain A sama dengan
kekayaan awalnya. Gunakan persamaan ini untuk menentukan peluang pemain A
bangkrut.

\i\label{exer 12.2.11} Dalam masalah kebangkrutan penjudi, andaikan
$p = q = 1/2$.
\begin{enumerate}
\item Dengan menggunakan Persamaan~\ref{eq 12.2.2}, bersama fakta bahwa
$q_0 = 1$ dan $q_M = 0$, tunjukkan bahwa untuk $0 \le z \le M$,
$$q_z = {{M - z}\over M}\ .$$
\item Dalam Persamaan~\ref{eq 12.2.3}, biarkan $p \rightarrow 1/2$ (dan karena
$q = 1 - p$, maka $q \rightarrow 1/2$ pula). Tunjukkan bahwa pada limitnya,
$$q_z = {{M - z}\over M}\ .$$  \emx {Petunjuk}: Ganti $q$ dengan $1-p$, dan
gunakan aturan L'Hopital.
\end{enumerate}

%We may want to move the next exercise to an earlier chapter.
\i\label{exer 12.2.12} Di kasino Amerika, roda rolet memuat bilangan bulat dari
1 sampai 36, bersama 0 dan 00. Separuh bilangan tak nol berwarna merah, separuh
lainnya hitam, sedangkan 0 dan 00 berwarna hijau. Taruhan yang lazim dalam
permainan ini adalah memasang satu dolar pada merah. Jika angka merah muncul,
petaruh menerima kembali dolarnya dan juga memperoleh satu dolar lagi. Jika
angka hitam atau hijau muncul, ia kehilangan dolarnya.
\begin{enumerate}
\item Misalkan seseorang memulai dengan 40 dolar dan terus bertaruh pada merah
hingga kekayaannya mencapai 50 atau 0. Tentukan peluang kekayaannya mencapai
50 dolar.
\item Berapa banyak uang yang harus ia miliki sebagai modal awal agar
peluangnya memenangkan 10 dolar sebelum bangkrut adalah 95\%?
\item Seorang pemilik kasino pernah berpendapat bahwa kasino tanpa 0 dan 00
tetap dapat memperoleh laba karena sebagian pelanggan akan terus bermain hingga
kehilangan seluruh uang mereka. Menurut Anda, apakah kasino seperti itu akan tetap
bertahan? Pernyataan ini dirumuskan ulang secara mandiri; ungkapan narasumber tidak
diterjemahkan atau direproduksi.
\end{enumerate}

\end{LJSItem}





 
\section{Hukum Arksinus}\label{sec 12.3}\index{hukum arksinus}
\par
Dalam Latihan~\ref{sec 12.1}.\ref{exer 12.1.6}, distribusi waktu keseimbangan
terakhir dalam perjalanan acak simetris telah ditentukan. Jika kita gunakan
$\alpha_{2k, 2m}$ untuk menyatakan peluang bahwa perjalanan acak sepanjang
$2m$ mencapai keseimbangan terakhirnya pada waktu $2k$, maka
$$\alpha_{2k, 2m} = u_{2k}u_{2m-2k}\ .$$
Sekarang kita akan menunjukkan cara menghampiri distribusi nilai-nilai $\alpha$
dengan sebuah fungsi sederhana. Ingat bahwa
$$u_{2k} \sim {1\over{\sqrt {\pi k}}}\ .$$
Karena itu, ketika $k$ dan $m$ keduanya menuju $\infty$, kita mempunyai
$$\alpha_{2k, 2m} \sim {1\over{\pi \sqrt{k(m-k)}}}\ .$$
Ekspresi terakhir ini dapat ditulis sebagai
$${1\over{\pi m \sqrt{(k/m)(1 - k/m)}}}\ .$$
Jadi, jika kita definisikan
$$f(x) = {1\over{\pi \sqrt{x(1-x)}}}\ ,$$
untuk $0 < x < 1$, maka
$$\alpha_{2k, 2m} \approx {1\over m}f\biggl({k\over m}\biggr)\ .$$
Tanda $\approx$ digunakan karena kita tidak lagi mensyaratkan $k$ menjadi
besar. Artinya, kita dapat mengganti distribusi diskret $\alpha_{2k, 2m}$
dengan densitas kontinu $f(x)$ pada interval $[0, 1]$ dan memperoleh hampiran
yang baik. Secara khusus, jika $x$ adalah bilangan real tetap di antara 0 dan
1, maka
$$\sum_{k < xm}\alpha_{2k, 2m} \approx \int_0^x f(t)\,dt\ .$$
Ternyata $f(x)$ mempunyai antiturunan yang sederhana, sehingga kita dapat
menulis
$$\sum_{k < xm}\alpha_{2k, 2m} \approx {2\over \pi}\arcsin \sqrt x\ .$$  
\index{hukum arksinus}
Dari grafik fungsi terakhir ini dapat dilihat bahwa fungsi tersebut mempunyai
minimum pada $x = 1/2$ dan simetris terhadap titik itu. Seperti dicatat dalam
latihan, hal ini menyiratkan bahwa separuh perjalanan sepanjang $2m$ tidak
mencapai keseimbangan lagi setelah waktu $m$, suatu fakta yang mungkin tidak
akan terduga.
\par
Ternyata densitas arksinus muncul dalam jawaban bagi banyak pertanyaan lain
tentang perjalanan acak pada garis. Ingat bahwa dalam Bagian~\ref{sec 12.1},
perjalanan acak dapat dipandang sebagai garis poligonal yang menghubungkan
$(0,0)$ dengan $(m, S_m)$. Berdasarkan penafsiran ini, kita definisikan
$b_{2k, 2m}$ sebagai peluang bahwa suatu perjalanan acak sepanjang $2m$
mempunyai tepat $2k$ dari $2m$ ruas garis poligonalnya di atas sumbu-$t$.
\par
Peluang $b_{2k, 2m}$ sering ditafsirkan melalui permainan dua pemain. (Pembaca
tentu ingat permainan Heads or Tails dalam Contoh~\ref{exam 1.3}.) Pemain A
dikatakan unggul pada waktu $n$ jika perjalanan acak berada di atas sumbu-$t$
pada waktu tersebut, atau jika perjalanan acak berada pada sumbu-$t$ pada
waktu $n$ tetapi berada di atas sumbu-$t$ pada waktu $n-1$. (Pada waktu 0,
tidak ada pemain yang unggul.) Kita dapat menanyakan berapa banyak waktu unggul
bagi pemain A yang paling mungkin dalam permainan sepanjang $2m$. Kebanyakan
orang akan menjawab $m$. Namun, teorema berikut menyatakan bahwa $m$ justru
merupakan banyaknya waktu unggul bagi pemain A yang paling kecil kemungkinannya,
sedangkan banyaknya waktu unggul yang paling mungkin adalah 0 atau $2m$.
\begin{theorem}\label{thm 12.3.1}
Jika Peter dan Paul memainkan permainan Heads or Tails sepanjang $2m$, peluang
Peter unggul tepat $2k$ kali sama dengan
$$
\alpha_{2k, 2m}\ .
$$
%\end{theorem}
\proof
Untuk membuktikan teorema ini, kita perlu menunjukkan bahwa
\begin{equation}  
b_{2k, 2m} = \alpha_{2k, 2m}\ .
\label{eq 12.3.1}  
\end{equation}
Latihan~\ref{sec 12.1}.\ref{exer 12.1.7} menunjukkan bahwa
$b_{2m, 2m} = u_{2m}$ dan $b_{0, 2m} = u_{2m}$, sehingga kita hanya perlu
membuktikan bahwa Persamaan~\ref{eq 12.3.1} berlaku untuk
$1 \le k \le m-1$. Kita dapat memperoleh rekursi yang melibatkan nilai-nilai
$b$ dan $f$ (yang didefinisikan dalam Bagian~\ref{sec 12.1}) dengan menghitung
banyaknya lintasan sepanjang $2m$ yang mempunyai tepat $2k$ ruas di atas
sumbu-$t$, dengan $1 \le k \le m-1$. Untuk menghitung kumpulan lintasan ini,
kita andaikan bahwa pengembalian pertama terjadi pada waktu $2j$, dengan
$1 \le j \le m-1$. Ada dua kasus yang perlu dipertimbangkan. Selama $2j$ hasil
pertama, lintasan berada di atas sumbu-$t$ atau di bawah sumbu-$t$. Dalam kasus pertama,
lintasan harus mempunyai tepat $(2k - 2j)$ ruas garis di atas sumbu-$t$ di
antara $t = 2j$ dan $t = 2m$. Dalam kasus kedua, lintasan harus mempunyai tepat
$2k$ ruas garis di atas sumbu-$t$ di antara $t = 2j$ dan $t = 2m$.
\par Sekarang kita hitung banyaknya lintasan dari berbagai jenis yang dijelaskan
di atas. Banyaknya lintasan sepanjang $2j$ yang semua ruas garisnya berada di
atas sumbu-$t$ dan kembali ke titik asal untuk pertama kalinya pada waktu $2j$
adalah $(1/2)2^{2j}f_{2j}$. Nilai ini juga sama dengan banyaknya lintasan
sepanjang $2j$ yang semua ruas garisnya berada di bawah sumbu-$t$ dan kembali
ke titik asal untuk pertama kalinya pada waktu $2j$. Banyaknya lintasan
sepanjang $(2m - 2j)$ yang mempunyai tepat $(2k - 2j)$ ruas garis di atas
sumbu-$t$ adalah $b_{2k-2j, 2m-2j}$. Terakhir, banyaknya lintasan sepanjang
$(2m-2j)$ yang mempunyai tepat $2k$ ruas garis di atas sumbu-$t$ adalah
$b_{2k,2m-2j}$. Karena itu,
$$b_{2k,2m} = {1\over 2} \sum_{j = 1}^k f_{2j}b_{2k-2j, 2m-2j} + 
{1\over 2}\sum_{j = 1}^{m-k} f_{2j}b_{2k, 2m-2j}\ .$$
\par
Sekarang kita andaikan bahwa Persamaan~\ref{eq 12.3.1} benar untuk $m < n$.
Maka

\begin{eqnarray*}
b_{2k, 2n} &=& {1\over 2} \sum_{j = 1}^k f_{2j}\alpha_{2k-2j, 2m-2j} + 
{1\over 2}\sum_{j = 1}^{m-k} f_{2j}\alpha_{2k, 2m - 2j}\\
&=& {1\over 2}\sum_{j = 1}^k f_{2j}u_{2k-2j}u_{2m-2k} + 
{1\over 2}\sum_{j = 1}^{m-k} f_{2j}u_{2k}u_{2m - 2j - 2k}\\
&=& {1\over 2}u_{2m-2k}\sum_{j = 1}^k f_{2j}u_{2k - 2j} +
{1\over 2}u_{2k}\sum_{j = 1}^{m-k} f_{2j}u_{2m - 2j - 2k}\\
&=& {1\over 2}u_{2m - 2k}u_{2k} + {1\over 2}u_{2k}u_{2m - 2k}\ ,
\end{eqnarray*}

di mana kesamaan terakhir mengikuti Teorema~\ref{thm 12.1.2}. Jadi,
$$b_{2k, 2n} = \alpha_{2k, 2n}\ ,$$
yang menyelesaikan pembuktian.
\end{theorem}
%\putfig{4.5truein}{PSfig12-2}{The arc sine density and times in the lead.}{fig 12.2}
Kita ilustrasikan teorema di atas dengan menyimulasikan 10{,}000 permainan
Heads or Tails, yang masing-masing terdiri atas 40 lemparan. Distribusi
banyaknya waktu Peter unggul diberikan dalam Gambar~\ref{fig 12.2}, bersama
densitas arksinus.

\putfig{3.5truein}{PSfig12-2}{Waktu dalam posisi unggul.}{fig 12.2}

Kita akhiri bagian ini dengan menyatakan dua hasil lain yang memuat densitas
arksinus. Bukti hasil-hasil ini dapat ditemukan dalam Feller.\footnote{W.~Feller,
op.\ cit., pp. 93--94.}
\begin{theorem}\label{thm 12.3.2}
Misalkan $J$ adalah peubah acak yang, untuk suatu perjalanan acak sepanjang
$2m$, memberikan subskrip terkecil $j$ sedemikian sehingga
$S_{j} = S_{2m}$. (Berdasarkan pertimbangan paritas, subskrip $j$ tersebut
harus genap.) Misalkan $\gamma_{2k, 2m}$ adalah peluang bahwa $J = 2k$. Maka
$$\gamma_{2k, 2m} = \alpha_{2k, 2m}\ .$$
\end{theorem}
Teorema berikut menyatakan bahwa densitas arksinus dapat diterapkan pada
berbagai macam situasi. Fungsi distribusi kontinu $F(x)$ disebut
\emx {simetris} jika $F(x) = 1 - F(-x)$. (Jika $X$ adalah peubah acak kontinu
dengan fungsi distribusi simetris, maka untuk setiap bilangan real $x$ kita
mempunyai $P(X \le x) = P(X \ge -x)$.) Bayangkan suatu perjalanan acak
sepanjang $n$ yang setiap sukunya mempunyai distribusi $F(x)$, dengan $F$
kontinu dan simetris. Subskrip \emx {maksimum pertama}\index{maksimum pertama
dari perjalanan acak} perjalanan tersebut adalah subskrip tunggal $k$
sedemikian sehingga
$$S_k > S_0,\ \ldots,\ S_k > S_{k-1},\ S_k \ge S_{k+1},\ \ldots,\ S_k \ge S_n\ .$$
Kita definisikan peubah acak $K_n$ sebagai subskrip maksimum pertama. Sekarang
kita dapat menyatakan teorema berikut mengenai peubah acak $K_n$.
\begin{theorem}\label{thm 12.3.3}
Misalkan $F$ adalah fungsi distribusi kontinu yang simetris, dan misalkan
$\alpha$ adalah bilangan real tetap yang terletak secara ketat di antara 0 dan
1. Ketika $n \rightarrow \infty$, kita mempunyai
$$P(K_n < n\alpha) \rightarrow {2\over \pi} \arcsin\sqrt \alpha\ .$$
\end{theorem}
\par
Versi teorema ini yang berlaku untuk perjalanan acak simetris juga dapat
ditemukan dalam Feller.

\exercises
\begin{LJSItem}

\i\label{exer 12.3.1} Untuk perjalanan acak sepanjang $2m$, definisikan
$\epsilon_k$ sama dengan 1 jika $S_k > 0$, atau jika $S_{k-1} = 1$ dan
$S_k = 0$. Definisikan $\epsilon_k$ sama dengan -1 dalam semua kasus lainnya.
Jadi, $\epsilon_k$ menunjukkan sisi sumbu-$t$ tempat perjalanan acak berada
selama interval waktu $[k-1, k]$. Suatu ``hukum bilangan besar" untuk barisan
$\{\epsilon_k\}$ akan menyatakan bahwa untuk setiap $\delta > 0$ berlaku
$$P\biggl(-\delta < {{\epsilon_1 + \epsilon_2 + \cdots + \epsilon_n}\over{n}} < \delta \biggr)
\rightarrow 1$$
ketika $n \rightarrow \infty$. Meskipun nilai-nilai $\epsilon$ tidak saling
independen, pernyataan di atas tentu tampak masuk akal. Dengan menggunakan
Teorema~\ref{thm 12.3.1}, tunjukkan bahwa jika $-1 \le x \le 1$, maka
$$\lim_{n \rightarrow \infty} P\biggl({{\epsilon_1 + \epsilon_2 + \cdots + \epsilon_n}\over{n}}
< x\biggr) = {2\over{\pi}} \arcsin\sqrt{{{1 + x}\over{2}}}\ .$$
 
\i\label{exer 12.3.2} Diberikan perjalanan acak $W$ sepanjang $m$, dengan
suku-suku
\[ \{X_1, X_2, \ldots,X_m\}\ , \]  
\noindent definisikan perjalanan acak \emx {terbalik} sebagai perjalanan $W^*$
dengan suku-suku
\[ \{X_m, X_{m-1}, \ldots, X_1\}\ . \]  
\begin{enumerate}
\item Tunjukkan bahwa jumlah parsial ke-$k$, yaitu $S^*_k$, memenuhi persamaan
$$S^*_k = S_m - S_{n-k}\ ,$$
di mana $S_k$ adalah jumlah parsial ke-$k$ untuk perjalanan acak $W$.
\item Jelaskan hubungan geometris antara grafik suatu perjalanan acak dan
grafik perjalanan terbaliknya. (Secara umum, tidak benar bahwa salah satu
grafik diperoleh dari grafik lainnya dengan pencerminan terhadap garis
vertikal.)
\item Gunakan bagian (a) dan (b) untuk membuktikan
Teorema~\ref{thm 12.3.2}.
\end{enumerate}

\end{LJSItem}

 
